\documentclass[tikz,border=5pt]{standalone}
\usetikzlibrary{calc,positioning}

%\nopagecolor
%\pagecolor{black!90} 

\begin{document}
\begin{tikzpicture}[
    dim/.style={draw=white, thin},
    bitfont/.style={
        font=\ttfamily\normalsize,
        text=white,
        inner sep=2pt
    },
    dimlabel/.style={
        font=\scriptsize\sffamily,
        text=white,
        align=center,
        inner xsep=4pt,
        inner ysep=0pt
    }
]

% 48-bit section (Chunk)
\node[bitfont] (chunk) {0110 1011 0010 \dots\ 1101 0011 1010};

% 16-bit section (Element)
\node[bitfont, right=8pt of chunk] (element) {1101 \dots\ 0011};

% Bracket above Chunk Index (starts 2mm above the text)
\path (chunk.north west) -- coordinate[midway] (c_mid) (chunk.north east);
\node[dimlabel] (c_lbl) at ([yshift=5mm]c_mid) {Chunk Index \\ (48 bits)};
\draw[dim] ([yshift=2mm]chunk.north west) |- (c_lbl) ([yshift=2mm]chunk.north east) |- (c_lbl);

% Bracket above Element Index (starts 2mm above the text)
\path (element.north west) -- coordinate[midway] (e_mid) (element.north east);
\node[dimlabel] (e_lbl) at ([yshift=5mm]e_mid) {Element Index \\ (16 bits)};
\draw[dim] ([yshift=2mm]element.north west) |- (e_lbl) ([yshift=2mm]element.north east) |- (e_lbl);

% Bit boundary numbers below
\node[font=\tiny\ttfamily, text=white, below=1pt of chunk.south west]   {63};
\node[font=\tiny\ttfamily, text=white, below=1pt of chunk.south east]   {16};
\node[font=\tiny\ttfamily, text=white, below=1pt of element.south west] {15};
\node[font=\tiny\ttfamily, text=white, below=1pt of element.south east] {0};

\end{tikzpicture}
\end{document}
